\documentclass{article}
\usepackage{graphicx}
\usepackage{amsfonts}
\usepackage{amsmath}
\usepackage{amsthm}
\usepackage{tikz}

\theoremstyle{definition}
\newtheorem{theorem}{Theorem}[section]
\newtheorem{corollary}[theorem]{Corollary}
\newtheorem{proposition}[theorem]{Proposition}
\newtheorem{lemma}[theorem]{Lemma}
\newtheorem{definition}[theorem]{Definition}

\theoremstyle{remark}
\newtheorem{remark}[theorem]{Remark}
\newtheorem{example}[theorem]{Example}
\newtheorem{claim}[theorem]{Claim}

\title{Comparison Test}
\author{Sarah Liu}
\date{January 2026}

\begin{document}

\maketitle

\section{Improper Integral}
\subsection{$p$-Test}
\subsection{Basic Comparison Test}
\subsection{Limit Comparison Test}

\section{Series}
\subsection{Integral Test}

\begin{theorem}
    Let $f$ be a continuous, positive, decreasing function on $[1,\infty)$ and let $a_n=f(n)$.
    Then the series $\sum_{n=1}^\infty a_n$ is convergent if and only if the improper integral $\int_1^{\infty}f(x)dx$ is convergent.
\end{theorem}

\subsection{Alternating Series Test}

\begin{theorem}
    If the alternating series
    \[
    \sum_{n=1}^{\infty}(-1)^{b+1}b_n=b_1-b_2+b_3-b_4+\dots \qquad b_n>0
    \]
    satisfies
    \begin{enumerate}
        \item $b_{n+1}\leq b_n \quad \text{for all }n$
        \item $\lim_{n\rightarrow \infty}b_n =0$
    \end{enumerate}
    then the series is convergent.
\end{theorem}

\section{Absolute Convergent}

\begin{theorem}
    If a series is absolutely convergent, then it is convergent.
\end{theorem}

\subsection{Ratio Test}

\begin{theorem}
Let $L=\lim_{n\rightarrow \infty}|\frac{a_{n+1}}{a_n}|$.
    \begin{enumerate}
        \item If $L<1$, then the series is absolutely convergent
        \item If $L>1$ (include $\infty$), then the series is divergent
        \item If $L=1$, then no conclusion can be drawn from the series
    \end{enumerate}
\end{theorem}

\end{document}